%
% 04_eigenschaften.tex -- Eigenschaften
%
% (c) 2020 Prof Dr Andreas Müller, Hochschule Rapperswil
%
% !TEX root = ../../buch.tex
% !TEX encoding = UTF-8
%

\section{Eigenschaften und Vergleiche mit anderen Methoden\label{qa:section:eigenschaften}}
\kopfrechts{Eigenschaften}
\subsection{Mehrfache Multiplikationen}
Im letzten Kapitel wurde gezeigt, dass für eine verzerrungsfreie 
Drehung eines Punktes $P$, der Punkt zwischen $q$ und $q^{-1}$ 
eingeklemmt werden muss. Mit der Eigenschaft, dass zwei negative 
Multiplikationen sich gegenseitig aufheben, kann die  
Gleichung
\begin{equation*}
    q P q^{-1} = (-q) P (-q)^{-1}
\end{equation*}
aufgestellt werden. Dies zeigt, dass die Multiplikation mit $q$ und 
$-q$ die gleiche Drehung erzeugt. Je nachdem, ob $q$ oder $-q$ 
verwendet wird, können für die gleiche Berechung mit 
Computeralgebra Rechenoperationen gespart werden.

\subsection{Vergleich mit Drehmatrizen}
Die Tabelle \ref{qa:tab:vergleich} zeigt die 
Gegenüberstellung beider Methoden für die Beschreibung einer 
Drehung im 3D Raum.
\begin{table}
        \centering
        \small
        \begin{tabular}{llll}
            \textbf{Eigenschaft} & \textbf{Quaternionen} & \textbf{Drehmatrizen}\\
            \hline
            Anzahl Parameter        & 4 Komponenten & 9 Matrixelemente  \\
	    davon unabhängig        & 3             & 3             \\
            Normalisierung          & Länge 1       & orthogonal    \\
            Parameterabhängigkeiten & $|q| = 1$     & $O^tO=I$      \\
            Anzahl Abhängigkeiten   & 1             & 6             \\
            Inverse                 & konjugiert: $q^{-1}=q^*$   & transponiert: $O^{-1}=O^t$  \\
            \hline
        \end{tabular}
        \caption{Vergleich von Quaternionen und Drehmatrizen.
	Das Produkt zwei $O^tO$ ist eine symmetrische Matrix, somit
	sind 6 der 9 Matrixelemente unabhängig. Dies erklärt, warum
	die Gleichung $O^tO=I$ den Parametern 6 Abhängigkeiten
        auferlegt.}
        \label{qa:tab:vergleich}
\end{table}
Für Drehmatrizen gilt die Bedingung, dass sie orthogonal sein 
müssen. Das bedeutet, dass aus den 9 Parametern nur 3 frei gewählt 
werden können. Die restlichen 6 Parameter müssen so gewählt werden, 
dass die Orthogonalität der Matrix erhalten bleibt. Bei 
Quaternionen reicht die Bedingung, dass die Norm 1 betragen muss, 
und somit können 3 aus den 4 Parametern frei gewählt werden.

\subsection{Vergleich mit Eulerwinkeln}
Eulerwinkel sind eine weitere Möglichkeit, eine Drehung im 3D-Raum 
\index{Euler-Winkel}%
zu beschreiben \cite[pp.~511-516]{qa:linalg}. Sie bestehen aus drei Winkeln, welche die Drehung 
um die drei Achsen beschreiben. Die Eulerwinkel haben den Vorteil, 
dass sie leicht verständlich sind und intuitiv verwendet werden 
können. Der Nachteil ist jedoch, dass sie nicht eindeutig sind und 
es zum Gimbal Lock kommen kann, was bedeutet, dass zwei Achsen in 
\index{Gimbal Lock}%
einer Linie liegen und somit ein Freiheitsgrad verloren geht. 
Quaternionen haben diesen Nachteil nicht und sind daher für viele 
Anwendungen besser geeignet. In der Astronomie finden Eulerwinkel 
\index{Astronomie}%
jedoch noch weiterhin Anwendung, da sich damit die Lage einer 
Planetenbahn im Raum auf natürliche Weise beschreiben und berechnen
\index{Planetenbahn}%
lässt.
